\begin{abstract}
Many advanced reactor designs require fuel enriched 
between 5 and 20\% $^{235}$U. To assist in producing 
fuel at these enrichment levels, government-owned 
inventories of highly-enriched uranium can be downblended. However, 
fuel produced from these inventories contain uranium 
impurities that are not often found when enriching natural 
uranium or accounted for when modeling reactor cores. To 
address this concern, this work models the X-energy Xe-100 and 
USNC Micro Modular Reactor 
designs and compares their performance with  
fuel from enriching natural uranium to fuel from downblended 
highly enriched uranium. This paper compares the models based on 
the effective neutron 
multiplication factor (\keff), effective delayed 
neutron fraction (\betaEff), energy- and spatially dependent 
neutron flux ($\phi$), as well as the fuel, coolant, moderator, and total reactivity 
temperature feedback coefficients ($\alpha_F$, $\alpha_C$, 
$\alpha_M$, $\alpha_T$). The results 
show that the fuel from downblended highly enriched uranium 
stockpiles leads to differences in each of the metrics, such as
a 1400 pcm increase in \keff for the Xe-100 design and a 700 pcm 
decrease in \keff at the end of cycle for the Micro Modular 
Reactor design.
However, the differences are not always statistically significant, 
nor do they inhibit reaching key design 
parameters, such as the 20 year cycle length of the Micro Modular 
Reactor design. 
\end{abstract}
